\documentclass{article}
\usepackage{graphicx}
\usepackage{amsthm}
\usepackage{amssymb}
\usepackage{amsmath}
\usepackage{parskip}
\usepackage[margin=0.6in]{geometry}
\usepackage{tikz}
\usepackage{titling}

 % move title upward

\title{15.095 Machine Learning HW3}
\author{Chris Low}
\date{} % removes the date

\begin{document}

\maketitle

\section*{Question 1}

\subsection*{(a)}

\textbf{False.}

Even if every tree is allowed to consider all $p$ features at every split, the trees are still trained on different bootstrap samples of the training data. Since bootstrap sampling is done with replacement, some observations may appear multiple times while others may not appear at all. Therefore, different trees may choose different optimal splits and need not be identical. Thus, the forest's majority-vote prediction need not be the same as the prediction of a single CART tree.

For example, suppose the original training observations are $
\{1,2,3,4,5\}.
$ Two bootstrap samples could be
$
\{2,3,1,2,5\}$ and 
$
\{3,1,4,5,1\}.
$

The first sample omits observation $4$ and repeats observation $2$, while the second omits observation $2$ and repeats observation $1$. Hence, even when all $p$ features are available at every split, the two trees are trained on different data and may select different splits.

\vspace{5mm}

\subsection*{(b)}

\textbf{True. 
}

Since feature $j$ only takes positive values, $\log(\cdot)$ is strictly increasing. This means that taking the log changes the numerical values of the feature, but does not change their ordering. Therefore, for any threshold $b>0$,

$$
x_{ij} < b
\iff
\log x_{ij} < \log b.
$$

So if the original OCT has a split $x_{ij}<b$, we can replace it with $\log x_{ij}<\log b$ after the transformation. The exact same observations will still go to the left and right sides of the split.

This also works in the reverse direction. If the transformed tree has a split $\log x_{ij}<b'$, then we can use $
x_{ij}<e^{b'}
$ on the original data. Thus, any split on one dataset has an equivalent split on the other dataset.

As a result, we can construct the same tree on both datasets, with the same observations in each leaf, the same leaf sizes, the same number of splits, and the same number of misclassified observations. The OCT objective is

$$
\min \sum_{t\in\mathcal{L}} L_t
+
\alpha \sum_{t\in\mathcal{B}} d_t,
$$

where the first term represents misclassification error and the second term penalizes the number of splits. Since neither of these changes under the log transformation, the optimal objective value is the same for the original and transformed data.


\newpage
\subsection*{(c)}

\textbf{True.}

Since $z_{it}=1$ if observation $i$ is assigned to leaf $t$, the total number of training observations assigned to leaf $t$ is $
\sum_{i=1}^{n} z_{it}.
$

We want every leaf that predicts class $A$ to contain at least $N_A$ training observations. Since $c_{At}=1$ if leaf $t$ predicts class $A$, we can add the constraint

$$
\sum_{i=1}^{n} z_{it} \geq N_A c_{At},
\qquad \forall t \in \mathcal{L}.
$$

If leaf $t$ predicts class $A$, then $c_{At}=1$, and the constraint becomes $
\sum_{i=1}^{n} z_{it} \geq N_A,
$ so at least $N_A$ training observations must be assigned to that leaf. If leaf $t$ does not predict class $A$, then $c_{At}=0$, and the constraint becomes $
\sum_{i=1}^{n} z_{it} \geq 0,
$ which adds no additional restriction. Therefore, this requirement can be added using one linear constraint per leaf and no additional variables.


\vspace{5mm}
\subsection*{(d)}

\textbf{False.}

The CART prediction is the average demand of the observations in the leaf:

$$
\hat{y}
=
\frac{40+50+60+70+130}{5}
=
70.
$$

However, the weighted-cost prescription does not simply set $z=\hat{y}$. Since all five observations are in the same CART leaf, they each receive weight $1/5$. Therefore, the prescription solves

$$
\min_{z\geq 0}
\frac{1}{5}
\sum_{i=1}^{5}
\left(7z-10\min(z,y_i)\right).
$$

To find the optimal $z$, consider the effect of baking one additional loaf. If demand is greater than the number of loaves baked, $y>z$, then all loaves are sold and $
c(z;y)=7z-10z=-3z.
$ Thus, each additional loaf decreases cost by $3$ (equivalently, increases profit by $3$).

On the other hand, if demand is less than the number of loaves baked, $y<z$, then only $y$ loaves are sold and \newline $
c(z;y)=7z-10y.
$ Since $y$ is fixed, increasing $z$ by one increases cost by $7$. Therefore, an additional loaf contributes $-3$ to cost if it is sold and $+7$ if it is not sold.

For $z<40$, the additional loaf would be sold on all five historical days, so the slope of the total cost is

$$
5(-3)=-15<0.
$$

For $40<z<50$, the additional loaf would be unsold on one day and sold on four days, giving

$$
7+4(-3)=-5<0.
$$

For $50<z<60$, it would be unsold on two days and sold on three days, giving

$$
2(7)+3(-3)=5>0.
$$

Thus, the objective decreases up to $z=50$ and increases after $z=50$, so
$ z^*=50.
$

Therefore, although the CART tree predicts demand of $70$, the weighted-cost prescription bakes $50$ loaves, not $70$.


\newpage

\subsection*{(e)}

\textbf{False.}

Under inverse propensity weighting, the estimated reward for point $i$ under treatment $t$ is

$$
\Gamma_{it}
=
\mathbf{1}\{T_i=t\}\frac{y_i}{p_{it}}.
$$

Therefore, for the treatment that point $i$ actually received,

$$
\Gamma_{iT_i}
=
\frac{y_i}{p_{iT_i}},
$$

while for every other treatment $t\neq T_i$, we have $
\Gamma_{it}=0.
$

If each training point is placed in its own leaf, the policy chooses the treatment with the largest estimated reward. Since $y_i\in\mathbb{R}$, it is possible that $y_i<0$. Because $p_{iT_i}>0$, this would imply

$$
\Gamma_{iT_i}
=
\frac{y_i}{p_{iT_i}}
<0,
$$

while every treatment not received has reward $0$.

For example, suppose point $i$ received treatment $1$, with $
y_i=-5
$ and $
p_{i1}=0.5.
$ 

Then $
\Gamma_{i1}
=
\frac{-5}{0.5}
=
-10,
$ while for another treatment, say treatment $2$, $
\Gamma_{i2}=0.
$ 

Since the policy maximizes reward and $0>-10$, it would prescribe treatment $2$ rather than the treatment point $i$ actually received. Therefore, the tree does not necessarily prescribe every training point its observed treatment.



\newpage

\section*{Question 2}


\subsection*{(a)}

At the root, the four possible splits are $x_1<1$, $x_1<2$, $x_2<1$, and $x_2<2$. For each child node $t$, the Gini impurity is

$$
I_t=1-\sum_{k\in\{A,B\}}p_{kt}^2,
$$

and the weighted Gini of a split is

$$
\frac{n_L}{n}I_L+\frac{n_R}{n}I_R,
$$

where $n=31$.

For the split $x_1<1$, the left child contains $(A,B)=(6,4)$ and the right child contains $(A,B)=(13,8)$. Thus,

$$
I_L
=
1-\left(\frac{6}{10}\right)^2-\left(\frac{4}{10}\right)^2
=
0.48,
$$

$$
I_R
=
1-\left(\frac{13}{21}\right)^2-\left(\frac{8}{21}\right)^2
\approx
0.4717,
$$

$$
G(x_1<1)
=
\frac{10}{31}(0.48)
+
\frac{21}{31}(0.4717)
\approx
0.4743.
$$

For the split $x_1<2$, the left child contains $(A,B)=(10,7)$ and the right child contains $(A,B)=(9,5)$. Thus,

$$
I_L
=
1-\left(\frac{10}{17}\right)^2-\left(\frac{7}{17}\right)^2
\approx
0.4844,
$$

$$
I_R
=
1-\left(\frac{9}{14}\right)^2-\left(\frac{5}{14}\right)^2
\approx
0.4592,
$$

$$
G(x_1<2)
=
\frac{17}{31}(0.4844)
+
\frac{14}{31}(0.4592)
\approx
0.4730.
$$

For the split $x_2<1$, the left child contains $(A,B)=(2,5)$ and the right child contains $(A,B)=(17,7)$. Thus,

$$
I_L
=
1-\left(\frac{2}{7}\right)^2-\left(\frac{5}{7}\right)^2
\approx
0.4082,
$$

$$
I_R
=
1-\left(\frac{17}{24}\right)^2-\left(\frac{7}{24}\right)^2
\approx
0.4132,
$$

$$
G(x_2<1)
=
\frac{7}{31}(0.4082)
+
\frac{24}{31}(0.4132)
\approx
0.4121.
$$

For the split $x_2<2$, the left child contains $(A,B)=(14,6)$ and the right child contains $(A,B)=(5,6)$. Thus,

$$
I_L
=
1-\left(\frac{14}{20}\right)^2-\left(\frac{6}{20}\right)^2
=
0.42,
$$

$$
I_R
=
1-\left(\frac{5}{11}\right)^2-\left(\frac{6}{11}\right)^2
\approx
0.4959,
$$

$$
G(x_2<2)
=
\frac{20}{31}(0.42)
+
\frac{11}{31}(0.4959)
\approx
0.4469.
$$

Therefore, the four weighted Gini impurities are

$$
\begin{aligned}
G(x_1<1) &\approx 0.4743,\\
G(x_1<2) &\approx 0.4730,\\
G(x_2<1) &\approx 0.4121,\\
G(x_2<2) &\approx 0.4469.
\end{aligned}
$$

Since CART chooses the split with the smallest weighted Gini impurity, the root split is $
\boxed{x_2<1}.
$


After choosing $x_2<1$ as the root split, we split each child separately.

\vspace{5mm} 
\textbf{Left child: $x_2<1$
}


This node contains $(A,B)=(2,5)$, so $n=7$. The two useful candidate splits are $x_1<1$ and $x_1<2$.

For $x_1<1$, the left child contains $(A,B)=(1,0)$ and the right child contains $(A,B)=(1,5)$. Thus,

$$
I_L=0,
$$

$$
I_R
=
1-\left(\frac{1}{6}\right)^2-\left(\frac{5}{6}\right)^2
\approx 0.2778.
$$

$$
G(x_1<1)
=
\frac{1}{7}(0)+\frac{6}{7}(0.2778)
\approx 0.2381.
$$

For $x_1<2$, the left child contains $(A,B)=(2,0)$ and the right child contains $(A,B)=(0,5)$. Both are pure, so

$$
I_L=I_R=0,
$$

$$
G(x_1<2)=0.
$$

$$
\boxed{x_1<2}
$$

for the left child.

The resulting leaves are $(A,B)=(2,0)$, which predicts class $A$, and $(A,B)=(0,5)$, which predicts class $B$. Both leaves have zero misclassification error.

\vspace{3mm}
\textbf{Right child: $x_2\geq 1$
}

This node contains $(A,B)=(17,7)$, so $n=24$. The useful candidate splits are $x_1<1$, $x_1<2$, and $x_2<2$.

For $x_1<1$, the two children contain $(A,B)=(5,4)$ and $(A,B)=(12,3)$. Thus,

$$
I_L
=
1-\left(\frac{5}{9}\right)^2-\left(\frac{4}{9}\right)^2
\approx 0.4938,
$$

$$
I_R
=
1-\left(\frac{12}{15}\right)^2-\left(\frac{3}{15}\right)^2
=
0.32.
$$


$$
G(x_1<1)
=
\frac{9}{24}(0.4938)
+
\frac{15}{24}(0.32)
\approx 0.3852.
$$

For $x_1<2$, the two children contain $(A,B)=(8,7)$ and $(A,B)=(9,0)$. Thus,

$$
I_L
=
1-\left(\frac{8}{15}\right)^2-\left(\frac{7}{15}\right)^2
\approx 0.4978,
$$

$$
I_R=0.
$$

$$
G(x_1<2)
=
\frac{15}{24}(0.4978)
+
\frac{9}{24}(0)
\approx 0.3111.
$$

For $x_2<2$, the two children contain $(A,B)=(12,1)$ and $(A,B)=(5,6)$. Thus,

$$
I_L
=
1-\left(\frac{12}{13}\right)^2-\left(\frac{1}{13}\right)^2
\approx 0.1420,
$$

$$
I_R
=
1-\left(\frac{5}{11}\right)^2-\left(\frac{6}{11}\right)^2
\approx 0.4959.
$$

$$
G(x_2<2)
=
\frac{13}{24}(0.1420)
+
\frac{11}{24}(0.4959)
\approx 0.3042.
$$

Comparing the three candidate splits,

$$
\begin{aligned}
G(x_1<1) &\approx 0.3852,\\
G(x_1<2) &\approx 0.3111,\\
G(x_2<2) &\approx 0.3042.
\end{aligned}
$$

Therefore, CART chooses
$
\boxed{x_2<2}
$ for the right child.

The leaf with $1\leq x_2<2$ contains $(A,B)=(12,1)$, so it predicts class $A$ and makes $1$ error. The leaf with $x_2\geq2$ contains $(A,B)=(5,6)$, so it predicts class $B$ and makes $5$ errors.

Thus, the four final leaves make $
0+0+1+5=6
$ total errors, so the training error of the greedy CART tree is
$
\boxed{6}.
$


\vspace{5mm}

\subsection*{(b)}

Instead of choosing the root greedily, we fix each possible root split and find the best second-level split for each of its two children. The minimum total errors obtained for the four possible root splits are

$$
\begin{aligned}
x_1<1 &: \quad 1+4=5,\\
x_1<2 &: \quad 1+0=1,\\
x_2<1 &: \quad 0+6=6,\\
x_2<2 &: \quad 3+0=3,
\end{aligned}
$$

where the two numbers in each sum are the minimum errors from the left and right subtrees, respectively.

The smallest total error is obtained by choosing $
\boxed{x_1<2}
$ as the root.

For the left branch, $x_1<2$, the best split is $x_2<2$. This gives $
x_2<2: (A,B)=(10,1),
$ which predicts class $A$ and makes one error, and $
x_2\geq2: (A,B)=(0,6),
$ which predicts class $B$ and makes no errors.

For the right branch, $x_1\geq2$, the best split is $x_2<1$. This gives $
x_2<1:\quad (A,B)=(0,5),
$ which predicts class $B$ and makes no errors, and $
x_2\geq1:\quad (A,B)=(9,0),
$ which predicts class $A$ and also makes no errors.

Therefore, the total error is $
\boxed{1}.
$ 

\newpage
From part (a), the weighted Gini impurities of the four possible root splits were

$$
\begin{aligned}
G(x_2<1)&=0.4121,\\
G(x_2<2)&=0.4469,\\
G(x_1<2)&=0.4730,\\
G(x_1<1)&=0.4743.
\end{aligned}
$$

Thus, $x_1<2$, which gives the globally best tree, was only the third-best root split according to Gini impurity.

If the root splits are instead evaluated by their immediate misclassification errors, they give

$$
\begin{aligned}
x_2<1 &: \quad 9,\\
x_2<2 &: \quad 11,\\
x_1<1 &: \quad 12,\\
x_1<2 &: \quad 12.
\end{aligned}
$$

Thus, $x_1<2$ is actually tied for the worst root split when only the immediate misclassification error is considered.

CART cannot find the optimal tree because it grows the tree greedily. At each step, it chooses the split that looks best immediately and does not consider how that split affects later splits. Although $x_1<2$ does not perform well at the root by itself, it creates two branches that can be separated very well by the next splits. Therefore, the locally best split does not necessarily lead to the globally best tree.

Finally, no tree of any depth using the allowed integer thresholds can achieve an error below $1$. The cell
$
\boxed{(0,1)\times(1,2)}
$ contains $5$ class-$A$ patients and $1$ class-$B$ patient. Since integer-threshold splits cannot separate observations within the same unit cell, at least one patient in this cell must always be misclassified. Since the tree above achieves error $1$, it is globally optimal.

\subsection*{(c)}

\subsubsection*{(i.)}

The objective is $
\operatorname{error}(T)
+
\alpha\cdot(\text{number of splits in }T).
$ From the previous parts, the best achievable classification errors using $0,1,2,$ and $3$ splits are
$$
E_0=12,\qquad
E_1=9,\qquad
E_2=6,\qquad
E_3=1.
$$

Therefore, adding the complexity penalty $\alpha$ for each split gives the four possible objective values

$$
\begin{aligned}
J_0(\alpha) &= 12,\\
J_1(\alpha) &= 9+\alpha,\\
J_2(\alpha) &= 6+2\alpha,\\
J_3(\alpha) &= 1+3\alpha.
\end{aligned}
$$

The three-split tree remains optimal until its objective reaches that of the tree with no splits. Setting these equal gives $
1+3\alpha=12,
$ 

so $
\alpha=\frac{11}{3}.
$

The one- and two-split trees do not become optimal before this point, and for $\alpha>11/3$ they are both more expensive than the no-split tree. Therefore, the optimal tree is

$$
\begin{cases}
0\leq\alpha<\dfrac{11}{3}:&
\text{the three-split tree from part (b), with objective }1+3\alpha,\\[2mm]
\alpha=\dfrac{11}{3}:&
\text{the three-split tree and the no-split tree tie, with objective }12,\\[2mm]
\alpha>\dfrac{11}{3}:&
\text{the no-split tree, with objective }12.
\end{cases}
$$

Thus, trees with $1$ or $2$ splits are never the unique optimum (in fact, they are never optimal for any $\alpha$).

\newpage

\subsubsection*{(ii.)}

We now start with the greedy CART tree from part (a) and apply cost-complexity pruning. The full CART tree has three splits and a total classification error of $6$, so its objective is

$$
J_3^{\mathrm{CART}}(\alpha)=6+3\alpha.
$$

For the two-split tree, the best choice is to prune the right child split $x_2<2$. The two leaves on this branch, $(A,B)=(12,1)$ and $(A,B)=(5,6)$, are merged into one leaf containing $
(A,B)=(17,7).
$

This leaf predicts class $A$ and therefore makes $7$ errors. The left subtree still has zero errors, so the best two-split pruned tree has total error $7$ and objective

$$
J_2^{\mathrm{CART}}(\alpha)=7+2\alpha.
$$

Pruning again leaves only the root split $x_2<1$. The left leaf contains $(A,B)=(2,5)$ and therefore makes $2$ errors, while the right leaf contains $(A,B)=(17,7)$ and makes $7$ errors. Hence, the one-split tree has $
2+7=9
$ errors and objective

$$
J_1^{\mathrm{CART}}(\alpha)=9+\alpha.
$$

Finally, pruning the root gives a tree with no splits. Since the full dataset contains $19$ class-$A$ and $12$ class-$B$ observations, the tree predicts $A$ for everyone and makes $12$ errors. Thus,

$$
J_0^{\mathrm{CART}}(\alpha)=12.
$$

Therefore, the four objectives available through pruning are

$$
\begin{aligned}
J_3^{\mathrm{CART}}(\alpha)&=6+3\alpha,\\
J_2^{\mathrm{CART}}(\alpha)&=7+2\alpha,\\
J_1^{\mathrm{CART}}(\alpha)&=9+\alpha,\\
J_0^{\mathrm{CART}}(\alpha)&=12.
\end{aligned}
$$

The first transition occurs when $
6+3\alpha=7+2\alpha,
$ which gives $
\alpha=1.
$

The next transition occurs when $
7+2\alpha=9+\alpha,
$ which gives $
\alpha=2.
$

Finally, $
9+\alpha=12
$ gives $
\alpha=3.
$

Therefore, CART returns

$$
\begin{cases}
0\leq\alpha<1 &: \text{the full three-split CART tree},\\
\alpha=1 &: \text{the three- and two-split trees tie},\\
1<\alpha<2 &: \text{the two-split pruned tree},\\
\alpha=2 &: \text{the two- and one-split trees tie},\\
2<\alpha<3 &: \text{the one-split tree},\\
\alpha=3 &: \text{the one- and zero-split trees tie},\\
\alpha>3 &: \text{the zero-split tree}.
\end{cases}
$$

From part (i), the globally optimal objective is

$$
J^*(\alpha)
=
\begin{cases}
1+3\alpha, & 0\leq\alpha\leq\frac{11}{3},\\
12, & \alpha\geq\frac{11}{3}.
\end{cases}
$$

For $\alpha<11/3$, the globally optimal tree has a strictly smaller objective than the tree returned by CART. Once $\alpha\geq11/3$, the globally optimal solution has no splits and objective $12$, which CART can also obtain by pruning the entire tree.

Therefore, CART achieves the same objective value as the globally optimal tree exactly when

$$
\boxed{\alpha\geq\frac{11}{3}}.
$$




\subsection*{(d)}

\subsubsection*{(i.)}

For each feature $j\in[p]$ and leaf $t\in\mathcal{L}$, introduce a continuous variable

$$
u_{jt}\in[0,1],
$$

where $u_{jt}$ indicates whether feature $j$ is tested at least once on the path from the root to leaf $t$.

For example, suppose the path to leaf $t$ contains two nodes, and both nodes split on feature $x_2$. Then

$$
a_{2,1}=1
\qquad\text{and}\qquad
a_{2,2}=1,
$$

since $a_{jm}$ indicates whether feature $j$ is used at a specific node $m$. However,

$$
u_{2t}=1,
$$

since $u_{jt}$ only indicates whether feature $j$ is used at least once anywhere on the path to leaf $t$. Thus, even though $x_2$ is used twice, its cost is counted only once.

For every ancestor $m\in P_t$, we impose

$$
u_{jt}\geq a_{jm},
\qquad
\forall j\in[p],\ t\in\mathcal{L},\ m\in P_t.
$$

Thus, if any node on the path to leaf $t$ splits on feature $j$, then $u_{jt}$ is forced to $1$. The budget constraint for each leaf is then

$$
\sum_{j=1}^p c_j u_{jt}\leq C,
\qquad
\forall t\in\mathcal{L}.
$$


\vspace{5mm}
There are $p|\mathcal{L}|$ additional continuous variables because for each of the $|\mathcal{L}|$ leaves, we need one variable for each of the $p$ possible features. For example, if $p=2$ and there are four leaves, we need

$$
u_{1,t_1},u_{2,t_1},
u_{1,t_2},u_{2,t_2},
u_{1,t_3},u_{2,t_3},
u_{1,t_4},u_{2,t_4},
$$

for a total of $
2(4)=8
$ variables. These variables record whether each feature is used on the path to each leaf. They do not imply that every feature is actually used.

Each leaf has at most $D$ ancestor nodes. Therefore, for each feature-leaf pair $(j,t)$, there are at most $D$ linking constraints of the form $
u_{jt}\geq a_{jm}.
$ Since there are $p|\mathcal{L}|$ feature-leaf pairs, there are at most $
pD|\mathcal{L}|
$ linking constraints.

In addition, there is one budget constraint for each leaf, giving another $
|\mathcal{L}|
$ constraints. Therefore, the total number of additional constraints is at most

$$
pD|\mathcal{L}|+|\mathcal{L}|
=
\boxed{|\mathcal{L}|(pD+1).}
$$

\newpage

\subsubsection*{(ii.)}


We have $$ 
c_1=1,\qquad c_2=3,\qquad C=3.
$$

Since using both features on the same path would cost $
c_1+c_2=1+3=4>C,
$ each path can use only $x_1$ or only $x_2$. An $x_1$-only tree has minimum error $12$ because no matter how we split by $x_1$ ($x_1 < 1$ or $x_1 < 2$), we would get an $A$ majority. Using only $x_2$, we can split first at $x_2<1$ and then split the right child at $x_2<2$. The three leaves have class counts

$$
(2,5),\qquad (12,1),\qquad (5,6),
$$

so their respective errors are $2$, $1$, and $5$. Thus, the total error is $
2+1+5=8.
$ Therefore, the minimum error under the correct budget constraint is $
\boxed{8}.
$

Under the colleague's constraint,
$$
\sum_{m\in P_t}\sum_{j=1}^p c_j a_{jm}\leq C,
$$

a feature is charged every time it is tested. Thus, a path that tests $x_2$ twice is charged $
3+3=6>C,
$ so the two-split $x_2$ tree above is no longer feasible. The best feasible tree instead uses the single split $x_2<1$. Its two leaves contain $(2,5)$ and $(17,7)$ observations, giving $
2+7=9
$ errors. Hence, the smallest error under the colleague's constraint is $
\boxed{9}.
$

The two answers differ because the colleague's constraint charges for every use of a feature, while the correct formulation charges each feature only once along a path. Thus, the colleague's constraint is too restrictive and rules out the $x_2$-only tree with error $8$.

\end{document}
